\documentclass[11pt]{article}
\usepackage{tcee}
\addbibresource{3B29.bib}
\tcetema[Financiación exterior, deuda externa y ayuda al desarrollo]{3.B.29}{La financiación exterior del desarrollo económico.\cortetitulo El problema de la deuda externa. La ayuda al desarrollo.}
\tcefecha{12/11/2026}
\begin{document}
\tceetitulo

\section*{Introducción}
\minutos{2,5}

\begin{itemize}
\item \textbf{\emph{Enganche}}:
  \begin{itemize}
  \item «Dale un pez a un hombre y comerá un día; enséñale a pescar y comerá toda la vida», dice el refrán.
  \item Del 30 de junio al 3 de julio de 2025, Sevilla acogió la IV Conferencia Internacional sobre la Financiación para el Desarrollo. Su documento final, el \emph{Compromiso de Sevilla}, parte de una cifra: a los países en desarrollo les faltan unos 4 billones de dólares al año para cumplir los Objetivos de Desarrollo Sostenible\nocite{ONU2025Sevilla,ONU2025FfD4Cierre}.
  \item Ese mismo año, la ayuda oficial de los donantes de la OCDE cayó un 23,1~\% en términos reales, la mayor caída registrada, hasta el 0,26~\% de su renta nacional bruta\nocite{OCDE2026AODPreliminar}, mientras los países de renta baja y media pagaban 1,4 billones de dólares por el servicio de su deuda externa\nocite{BancoMundial2026IDR}.
  \item La pregunta que recorre la exposición es: \emph{¿por qué el capital que necesitan los países pobres no llega, llega caro o llega regalado y no sirve?}
  \end{itemize}

\item \textbf{\emph{Relevancia}}:
  \begin{itemize}
  \item \emph{Para los países en desarrollo}: sin ahorro exterior, su inversión queda limitada por su propio ahorro, que es escaso precisamente porque son pobres; con demasiado, la deuda se come lo que crecen. Entre 2022 y 2025 pagaron 436.000 millones de dólares más por su deuda de lo que recibieron en financiación nueva\nocite{BancoMundial2026IDR}.
  \item \emph{Para la economía mundial}: las crisis de la deuda de los 80 y de los países más pobres en los 90 marcaron el sistema financiero internacional, y la ola de deuda posterior a la pandemia sigue abierta \vertema{3.B.21}.
  \item \emph{Para España y para el cuerpo}: España es miembro del Club de París, accionista de los bancos multilaterales y donante del Comité de Ayuda al Desarrollo; los técnicos comerciales y economistas del Estado la representan en esos foros y gestionan los instrumentos financieros de apoyo a la internacionalización y la conversión de deuda \vertema{4.A.24} \vertema{4.A.26}.
  \end{itemize}

\item \textbf{\emph{Contextualización}}:
  \begin{itemize}
  \item \emph{Las definiciones de partida}:
    \begin{itemize}
    \item los \concepto{países en desarrollo} son, en las estadísticas del Banco Mundial, los de \emph{renta baja y media}, clasificados por su renta nacional bruta por habitante; el desarrollo es más que crecer, pero sin crecer no hay desarrollo \vertema{3.B.28};
    \item la \concepto{financiación exterior} es el ahorro que el resto del mundo pone a disposición de un país: su contrapartida es un déficit por cuenta corriente, porque $\mathrm{CC}=S-I$, y el país invierte más de lo que ahorra \vertema{3.B.11};
    \item la \concepto{deuda externa bruta} es el saldo de los pasivos frente a no residentes que obligan a pagar principal o intereses (préstamos, bonos, depósitos, créditos comerciales); excluye las acciones y la inversión directa en capital\nocite{FMI2014DeudaExterna}. Restando los activos de deuda frente al exterior se obtiene la \emph{deuda externa neta};
    \item la \concepto{ayuda oficial al desarrollo} (AOD) es la parte de la financiación exterior que dan los gobiernos en condiciones de favor y con el desarrollo como fin principal (\ref{sec:aod-concepto}).
    \end{itemize}
  \item \emph{La historia}:
    \begin{itemize}
    \item la ayuda nace tras la Segunda Guerra Mundial, con el Plan Marshall y la descolonización, y se institucionaliza en 1961 con el Comité de Ayuda al Desarrollo (CAD) de la OCDE;
    \item la deuda vive 3 grandes episodios: la crisis latinoamericana de 1982, la deuda de los países más pobres, que acaba en las condonaciones de 1996-2005, y la nueva ola de deuda de los 2010, con China como gran acreedor y la pandemia como detonante;
    \item la agenda pasa de la ayuda a la financiación del desarrollo en su conjunto, con 4 conferencias de las Naciones Unidas: Monterrey (2002), Doha (2008), Addis Abeba (2015) y Sevilla (2025).
    \end{itemize}
  \end{itemize}

\item \textbf{\emph{Problemática}}:
  \begin{itemize}
  \item ¿Por qué debería el capital ir de los países ricos a los pobres, y por qué no lo hace?
  \item ¿Cuándo es sostenible la deuda externa, por qué pagan los Estados lo que deben y cómo se han resuelto las crisis de la deuda?
  \item ¿Qué es la ayuda al desarrollo, cómo ha evolucionado y funciona de verdad?
  \end{itemize}

\item \textbf{\emph{Estructura}}:
  \begin{itemize}
  \item La exposición sigue el título, de lo general a lo particular: de todas las fuentes de financiación exterior a la que crea deuda y, por último, a la que no se devuelve.
    \begin{itemize}
    \item \concepto{I. La financiación exterior del desarrollo} (\ref{sec:financiacion}): por qué financiarse fuera, por qué el capital no va a los pobres y qué fuentes hay;
    \item \concepto{II. El problema de la deuda externa} (\ref{sec:deuda}): sostenibilidad, impago y sobreendeudamiento, las crisis y sus respuestas y el marco institucional;
    \item \concepto{III. La ayuda al desarrollo} (\ref{sec:ayuda}): concepto y medida, evolución y agenda, y el debate sobre su eficacia.
    \end{itemize}
  \end{itemize}
\end{itemize}

\begin{esquema}[unbreakable]
\textbf{3.B.29. La financiación exterior del desarrollo económico. El problema de la deuda externa. La ayuda al desarrollo}\\[2pt]
Introducción\esqtiempo{2,5}\\
I. La financiación exterior del desarrollo\esqtiempo{8,5}\\
\hspace*{1.5em}1. ¿Por qué financiarse fuera?\\
\hspace*{1.5em}2. El capital no va a los pobres\\
\hspace*{1.5em}3. Las fuentes de la financiación exterior\\
II. El problema de la deuda externa\esqtiempo{9}\\
\hspace*{1.5em}1. Sostenibilidad, impago y sobreendeudamiento\\
\hspace*{1.5em}2. Las crisis de la deuda y sus respuestas\\
\hspace*{1.5em}3. El marco institucional y las soluciones\\
III. La ayuda al desarrollo\esqtiempo{8,5}\\
\hspace*{1.5em}1. Concepto y medida\\
\hspace*{1.5em}2. Evolución y agenda\\
\hspace*{1.5em}3. ¿Funciona la ayuda?\\
Conclusión\esqtiempo{1,5}
\end{esquema}

\begin{notaopositor}
Reparto del cante: introducción, \tiempo{2,5}; I, \tiempo{8,5} (I.1, \tiempo{3,5}; I.2, \tiempo{1,5}; I.3, \tiempo{3,5}); II, \tiempo{9} (II.1, \tiempo{3,5}; II.2, \tiempo{3,5}; II.3, \tiempo{2}); III, \tiempo{8,5} (III.1, \tiempo{2,5}; III.2, \tiempo{2,5}; III.3, \tiempo{3,5}), y conclusión, \tiempo{1,5}. Los programas del FMI y la estructura del Banco Mundial tienen su propio tema: aquí basta con nombrarlos. Los 3 gráficos imprescindibles son el del modelo intertemporal, el de la dinámica de la deuda externa y la curva de Laffer de la deuda; las funciones clave, las 2 brechas, la dinámica de la deuda externa y el saldo que la estabiliza. El refrán del enganche vuelve en la idea final. Las 3 cifras que conviene llevar al día son la AOD del CAD y su porcentaje de la renta, la deuda externa y su servicio en los países de renta baja y media y el número de países en riesgo alto de sobreendeudamiento.
\end{notaopositor}


% ==========================================================================
\section{La financiación exterior del desarrollo}\label{sec:financiacion}
\minutos{8,5}

\begin{itemize}
\item Empezamos por la financiación exterior en su conjunto: por qué la necesita un país pobre, por qué le llega tan poca y en qué formas le llega.
\item \emph{Conclusión anticipada}: la teoría dice que un país con poco capital debería recibirlo del resto del mundo y ganar con ello; los datos dicen que el capital va poco a los pobres y que importa tanto \emph{cómo} llega como \emph{cuánto} llega.
\end{itemize}

% --------------------------------------------------------------------------
\subsection{¿Por qué financiarse fuera?}\label{sec:brecha}

\begin{itemize}
\item Hay 2 respuestas: la de la economía del desarrollo de los 50 y 60, que calcula cuánto capital falta para crecer (la brecha de financiación), y la de la macroeconomía abierta, que pregunta cuándo le conviene a un país endeudarse (el enfoque intertemporal).
\end{itemize}

\subsubsection{La brecha de financiación}\label{sec:dos-brechas}

\paragraph{a) Idea}

\begin{itemize}
\item \emph{El punto de partida}: un país pobre ahorra poco porque es pobre, y sigue siendo pobre porque invierte poco. Es el \concepto{círculo vicioso de la pobreza} de \textcite{Nurkse1953}.
\item \emph{La salida}: si el ahorro interno no alcanza para la inversión que exige crecer, la diferencia tiene que ponerla el resto del mundo, con crédito o con ayuda. Es la justificación más antigua de la financiación exterior del desarrollo y la que dio origen a la ayuda.
\end{itemize}

\paragraph{b) Modelo}

\subparagraph{i) Supuestos}

\begin{itemize}
\item \emph{Tecnología}: la producción es proporcional al capital, con una relación capital-producto $v\equiv K/Y$ fija y sin sustitución entre factores. Es la versión de \autor{Evsey Domar} del modelo de \autor{Roy Harrod} y \autor{Evsey Domar} \vertema{3.A.43}.
\item \emph{Ahorro}: la tasa de ahorro interna, $s$, es exógena, y no hay depreciación.
\item \emph{Objetivo}: el Gobierno fija la tasa de crecimiento que quiere alcanzar, $\bar\gamma_Y$ (la barra marca lo fijado).
\end{itemize}

\subparagraph{ii) Desarrollo}

\begin{itemize}
\item \emph{La inversión necesaria}: como $Y=K/v$ e $I=\Delta K$, crecer a $\bar\gamma_Y$ exige invertir una parte de la renta igual a la relación capital-producto por el crecimiento objetivo, y lo que no cubre el ahorro interno es la \concepto{brecha de ahorro}:
  \begin{equation}\label{eq:brecha-ahorro}
  \frac{I}{Y}=v\,\bar\gamma_Y ,\qquad \frac{I-S}{Y}=v\,\bar\gamma_Y-s .
  \end{equation}
\item \emph{Su lectura}: $v\,\bar\gamma_Y-s$ es la financiación exterior «necesaria» para alcanzar el objetivo \vertema{3.B.28}.
  \begin{itemize}
  \item Con $v=3$ y un objetivo del 6~\%, hay que invertir el 18~\% del PIB; si el país ahorra el 12~\%, necesita financiación exterior por el 6~\% del PIB cada año.
  \end{itemize}
\end{itemize}

\subparagraph{iii) Implicaciones}

\begin{itemize}
\item \emph{La financiación exterior acelera el crecimiento en proporción a la brecha que cubre}: cada punto de PIB de crédito o de ayuda que se invierte añade $1/v$ puntos de crecimiento.
\item \emph{Una regla de cálculo}: durante décadas, el Banco Mundial y otros organismos calcularon con ella cuánta ayuda necesitaba cada país\nocite{Easterly1999}.
\end{itemize}

\paragraph{c) Extensiones: las 2 brechas}

\begin{itemize}
\item \emph{El segundo límite}: \textcite{CheneryStrout1966} observan que un país pobre importa casi todos sus bienes de equipo, de modo que crecer no exige solo ahorro, sino también divisas.\footnote{\textcite{CheneryStrout1966} añaden un tercer límite, la \emph{capacidad de absorción}: la falta de técnicos, de gestores y de instituciones pone un techo a la inversión que un país puede aprovechar, aunque tenga ahorro y divisas. Es la primera forma de una idea que vuelve en el debate sobre la eficacia de la ayuda (\ref{sec:aod-eficacia}).}
  \begin{itemize}
  \item La \emph{brecha de ahorro} es la inversión que exige el objetivo menos el ahorro interno, $I-S$.
  \item La \emph{brecha de divisas} son las importaciones que exige el objetivo menos lo que el país ingresa por exportaciones, $M-X$.
  \end{itemize}
\item \emph{La financiación exterior necesaria}, $\mathrm{FE}$, es la mayor de las 2 brechas ($\mathrm{FE}$ va en redonda porque $F$ es la función de producción):
  \funcionclave[eq:dos-brechas]{\frac{\mathrm{FE}}{Y}=\max\left\{v\,\bar\gamma_Y-s,\ \frac{M-X}{Y}\right\}}
\item \emph{Por qué el máximo}: después de los hechos las 2 brechas son iguales, porque el ahorro menos la inversión es el saldo exterior \vertema{3.B.11}; antes, en los planes, pueden no serlo, y la mayor es la que frena el crecimiento. La otra se cierra sola, con ahorro o con divisas que quedan sin usar.
\item \emph{Si manda la brecha de divisas}, el ahorro interno no basta aunque exista: no se puede convertir en las máquinas que hay que comprar fuera. Cada dólar de financiación exterior rinde entonces más que si solo faltara ahorro.
\end{itemize}

\paragraph{d) Evidencia empírica}

\begin{itemize}
\item \emph{La regla no funcionó}: la ayuda no se tradujo en inversión ni la inversión en crecimiento. Si toda la ayuda recibida por Zambia desde los 60 se hubiera invertido y la inversión se hubiera convertido en crecimiento como dice la regla, su renta por habitante habría superado los 20.000 dólares en 1994; era de unos 600\nocite{Easterly1999}.
\item \emph{El propio \autor{Evsey Domar}} había renegado de ese uso de su modelo, que había pensado para el ciclo y no para el crecimiento a largo plazo\nocite{Easterly2003}.
\end{itemize}

\paragraph{e) Valoración}

\begin{itemize}
\item \emph{Sus supuestos}: con rendimientos decrecientes y sustitución entre factores, acumular capital no sostiene por sí solo el crecimiento \vertema{3.A.43}; y la ayuda es fungible: puede sustituir al ahorro interno en lugar de sumarse a él (\ref{sec:aod-eficacia}).
\item \emph{Lo que no dice}: calcula cuánto falta, pero no si al país le conviene endeudarse ni cuánto. Para eso hace falta un modelo con optimización: el enfoque intertemporal.
\end{itemize}

\subsubsection{El enfoque intertemporal}\label{sec:intertemporal}

\paragraph{a) Idea}

\begin{itemize}
\item \emph{La pregunta}: ¿por qué le conviene a un país pobre endeudarse con el exterior?
\item \emph{La respuesta}: un déficit exterior no es malo en sí; es la forma en que un país pide prestado hoy para invertir o consumir y devuelve mañana. El enfoque intertemporal de la cuenta corriente lo formaliza y permite juzgar los déficits por su efecto sobre el bienestar\nocite{ObstfeldRogoff1996Foundations} \vertema{3.B.12}.
\end{itemize}

\paragraph{b) Modelo}

\subparagraph{i) Supuestos}

\begin{itemize}
\item \emph{Economía}: pequeña y abierta, con 2 periodos ($t$ y $t+1$), un solo bien que se consume o se invierte y previsión perfecta.
\item \emph{Preferencias}: un agente representativo con utilidad $U=u(C_t)+\beta\,u(C_{t+1})$, $\beta\equiv1/(1+\rho)$ \vertema{3.A.29}.
\item \emph{Tecnología}: la producción depende del capital, $F(K)$, con rendimientos decrecientes y sin depreciación ($K_{t+1}=K_t+I_t$).
\item \emph{Mercados}: acceso perfecto al mercado internacional de capitales al tipo real $r^*$, dado; todo déficit se financia con deuda.
\end{itemize}

\subparagraph{ii) Desarrollo}

\begin{itemize}
\item \emph{En autarquía}, el país solo puede pasar recursos al futuro invirtiendo. El equilibrio está donde la curva de indiferencia es tangente a la frontera de posibilidades de producción intertemporal (punto $A$ del gráfico~\ref{fig:cc-intertemporal}), y su pendiente común define el \concepto{tipo de interés de autarquía}, $-(1+r^A)$.
  \begin{itemize}
  \item Un país pobre tiene poco capital, que rinde mucho en el margen: su tipo de autarquía es mayor que el mundial, $r^*<r^A$.
  \end{itemize}
\item \emph{Con apertura}, la decisión se separa en 2, como en el \concepto{teorema de separación} de \autor{Irving Fisher} \vertema{3.A.29}:
  \begin{itemize}
  \item primero se invierte hasta que el capital rinde lo que cuesta pedir prestado fuera, sin mirar las preferencias:
    \begin{equation}\label{eq:inversion-abierta}
    F'(K_{t+1})=r^* ;
    \end{equation}
  \item después se reparte el consumo entre los 2 periodos, pidiendo prestado o prestando al tipo mundial.
  \end{itemize}
\item \emph{La cuenta corriente} es el ahorro menos la inversión, y el superávit comercial de $t+1$ devuelve lo pedido en $t$ con sus intereses:
  \begin{equation}\label{eq:cc-intertemporal}
  \mathrm{CC}_t=S_t-I_t ,\qquad \mathrm{XN}_{t+1}=-(1+r^*)\,\mathrm{CC}_t .
  \end{equation}
\item \emph{Interpretación gráfica} (gráfico~\ref{fig:cc-intertemporal}), con $r^*<r^A$:
  \begin{itemize}
  \item el país \emph{produce} en $P$, donde la frontera tiene la pendiente del mercado, $-(1+r^*)$: invierte más que en autarquía, porque el capital rinde dentro más de lo que cuesta fuera ($P$ es un punto, no el nivel de precios);
  \item \emph{consume} en $C$, en una curva de indiferencia más alta que la de $A$;
  \item en $t$ tiene un \emph{déficit por cuenta corriente}, la distancia horizontal entre $C$ y $P$, que se reparte entre más inversión ($\Delta I$) y más consumo ($\Delta C$) que en autarquía y se financia con una entrada neta de capital;
  \item en $t+1$ tiene un \emph{superávit comercial}, la distancia vertical entre $P$ y $C$, con el que devuelve la deuda y sus intereses.
  \end{itemize}
\end{itemize}

\begin{figure}[H]
\centering
\caption{El modelo intertemporal de la cuenta corriente: el país pobre que se abre}\label{fig:cc-intertemporal}
\grafico[0.66\textwidth]{cc-intertemporal}
\fuente{Elaboración propia a partir de \textcite[cap.~1]{ObstfeldRogoff1996Foundations}. $A$ es la autarquía; con $r^*<r^A$, el país produce en $P$, consume en $C$ y tiene un déficit en $t$ que se reparte entre más inversión ($\Delta I$) y más consumo ($\Delta C$).}
\end{figure}

\subparagraph{iii) Implicaciones}

\begin{itemize}
\item \emph{La dirección del capital}: \uline{el capital debería ir de los países ricos a los pobres, y los 2 ganarían}. El pobre invierte más y reparte mejor su consumo en el tiempo; el rico presta a un tipo mayor que el que obtendría en casa.
\item \emph{El bienestar}: la apertura nunca lo reduce, y corregir un déficit óptimo lo empeora.
\item \emph{El uso de la deuda}: un déficit que financia inversión rentable se paga solo; uno que financia consumo frente a una caída permanente de la renta solo retrasa el ajuste. Lo temporal se financia; lo permanente se ajusta\nocite{ObstfeldRogoff1996Foundations}.
\end{itemize}

\paragraph{c) Valoración}

\begin{itemize}
\item \emph{Sus supuestos}: acceso perfecto al crédito y ningún riesgo de impago, justo lo que menos se cumple en los países pobres; por qué pagan los Estados lo que deben es una pregunta del segundo bloque (\ref{sec:sostenibilidad-impago}).
\item \emph{Sus predicciones}: grandes flujos de capital de los países ricos a los pobres, que los datos no muestran. Es la paradoja de \autor{Lucas}.\footnote{\autor{Robert Lucas} recibió el Premio Nobel de Economía en 1995 «por haber desarrollado y aplicado la hipótesis de las expectativas racionales, transformando así el análisis macroeconómico y profundizando en la comprensión de la política económica»\nocite{NobelPrize1995Lucas}.}
\end{itemize}

% --------------------------------------------------------------------------
\subsection{El capital no va a los pobres}\label{sec:lucas}

\begin{itemize}
\item Una vez vista la teoría, pasamos a los hechos, que van en la dirección contraria.
\item \emph{La paradoja de \autor{Lucas}}: el capital no va de los países ricos a los pobres\nocite{Lucas1990Why}.
  \begin{itemize}
  \item Con la misma tecnología, una renta por persona 15 veces menor en la India que en Estados Unidos y una participación del capital de 0,4, el rendimiento del capital en la India debería ser 58 veces mayor.
  \item Con el capital humano la diferencia baja a 5, y con sus externalidades prácticamente desaparece.
  \item El capital va a donde ya hay capital, como en el evangelio: «al que tiene, se le dará» (Mateo 25, 29).
  \end{itemize}
\item \emph{Las explicaciones} son de 2 familias:
  \begin{enumerate}
  \item \emph{los fundamentales}: el capital rinde en los países pobres menos de lo que parece, porque tienen menos productividad total de los factores, menos capital humano, peores infraestructuras e instituciones más débiles, que según \autor{Daron Acemoglu}, \autor{Simon Johnson} y \autor{James Robinson} explican buena parte de las diferencias de renta entre países \vertema{3.B.28}\footnote{\autor{Daron Acemoglu}, \autor{Simon Johnson} y \autor{James Robinson} recibieron el Premio Nobel de Economía en 2024 «por sus estudios sobre cómo se forman las instituciones y cómo afectan a la prosperidad»\nocite{NobelPrize2024}.};
  \item \emph{los fallos del mercado internacional de capitales}:
    \begin{itemize}
    \item el riesgo de expropiación y de impago, porque no hay un tribunal que obligue a pagar a un Estado (\ref{sec:sostenibilidad-impago});
    \item la información asimétrica entre el deudor y el acreedor, que encarece el crédito al país pobre\nocite{GertlerRogoff1990};
    \item la escasez de activos financieros seguros en los países pobres, que concentra el ahorro mundial en los pocos países que los emiten\nocite{CaballeroFarhiGourinchas2017Safe}.
    \end{itemize}
  \end{enumerate}
\item \emph{La intolerancia a la deuda}: los países con un historial de impagos entran en crisis con deudas mucho menores que los demás, de modo que no hay un umbral de deuda que valga para todos\nocite{Reinhart2003Intolerancia}.
\item \emph{El puzle de la asignación}: entre los países en desarrollo, los que más crecen en productividad reciben menos capital neto, e incluso lo exportan, al revés de lo que predice el modelo; detrás están la acumulación de reservas y el ahorro público, no la inversión privada\nocite{GourinchasJeanne2013Capital} \vertema{3.B.12}.
\item \emph{Más capital exterior no es más crecimiento}: entre los países no industrializados, los que más se apoyaron en el ahorro exterior no crecieron más que los que se financiaron con el suyo, y la cuenta corriente y el crecimiento van juntos en sentido positivo\nocite{PrasadRajanSubramanian2007}.
\item \uline{El capital exterior complementa al ahorro interno, pero no lo sustituye}: importa la forma en que llega, que es lo que vemos a continuación.
\end{itemize}

% --------------------------------------------------------------------------
\subsection{Las fuentes de la financiación exterior}\label{sec:fuentes}

\begin{itemize}
\item Ya sabemos que la financiación exterior puede ser útil pero no llega sola; pasamos a sus formas.
\item La pregunta que las ordena es \emph{quién asume el riesgo}: si el país se compromete a pagar pase lo que pase (deuda) o si lo que paga depende de cómo le vaya (inversión directa, acciones).
\item \emph{Los órdenes de magnitud}: en 2025, los países de renta baja y media recibieron unos 700.000 millones de dólares en remesas, más que en inversión directa (unos 494.000 millones), y una entrada neta de deuda de 423.500 millones, más del doble que un año antes; la ayuda oficial que recibieron en 2023, último año con dato, fue de unos 252.000 millones\nocite{BancoMundial2026WDIFlujos,BancoMundial2026IDR} (gráfico~\ref{fig:flujos-desarrollo}).
\end{itemize}

\begin{figure}[H]
\centering
\caption{Los flujos de financiación exterior hacia los países de renta baja y media, 2000-2025}\label{fig:flujos-desarrollo}
\grafico[0.9\textwidth]{flujos-desarrollo}
\fuente{Elaboración propia a partir de \textcite{BancoMundial2026WDIFlujos} y \textcite{BancoMundial2026IDS}. Miles de millones de dólares corrientes. Inversión directa: entradas netas; remesas: recibidas por los hogares; AOD: ayuda neta recibida de todos los donantes y organismos (último dato, 2023); deuda: desembolsos menos amortizaciones de toda la deuda externa.}
\end{figure}

\subsubsection{Las que no crean deuda}\label{sec:sin-deuda}

\begin{itemize}
\item \emph{La inversión extranjera directa} (IED) da al inversor una influencia duradera en la empresa, con un 10~\% o más de su capital\nocite{FMI2009MBP6} \vertema{3.B.11}.
  \begin{itemize}
  \item \emph{Ventajas}: es la más estable, trae tecnología, gestión y acceso a mercados, y sus rentas caen en las malas épocas, de modo que el inversor comparte el riesgo.
  \item \emph{Por qué invierten fuera las empresas}: por sus ventajas de propiedad, de localización y de internalización, el paradigma ecléctico de \textcite{Dunning1980}; y solo las más productivas pueden pagar el coste fijo de producir fuera en lugar de exportar\nocite{HelpmanMelitzYeaple2004} \vertema{3.B.18}.
  \item \emph{Lo que controla el país} es la localización: estabilidad, infraestructuras, capital humano y reglas.
  \item \emph{Inconvenientes}: se concentra en pocos países de renta media y, en los más pobres, en los recursos naturales; y repatría sus beneficios.
  \end{itemize}
\item \emph{La inversión de cartera en acciones}:
  \begin{itemize}
  \item \emph{Ventajas}: no crea deuda y sus rentas caen en las crisis.
  \item \emph{Inconvenientes}: es volátil, con efecto rebaño, y exige mercados financieros profundos que los países pobres no tienen; en 2024 salió en términos netos de los países de renta baja y media\nocite{BancoMundial2026WDIFlujos}.
  \end{itemize}
\item \emph{Las remesas} de los emigrantes, hoy la mayor entrada de recursos de los países de renta baja y media:
  \begin{itemize}
  \item \emph{Ventajas}: son estables, a menudo anticíclicas (el emigrante envía más cuando su familia lo pasa mal) y llegan directamente a los hogares.
  \item \emph{Inconvenientes}: se dedican sobre todo al consumo, su efecto sobre el crecimiento es difícil de encontrar en los datos\nocite{ClemensMcKenzie2018}, pueden apreciar el tipo de cambio real y tienen detrás la salida de trabajadores, a menudo cualificados; además, enviarlas sigue costando más del 3~\% que fija como meta la Agenda 2030.
  \end{itemize}
\item \emph{Las donaciones oficiales}: la parte de la ayuda que no se devuelve (\ref{sec:ayuda}).
\end{itemize}

\subsubsection{Las que crean deuda}\label{sec:con-deuda}

\begin{itemize}
\item \emph{Los préstamos de bancos comerciales}: fueron la gran fuente de los 70, con el reciclaje de los petrodólares, y el origen de la crisis de 1982 (\ref{sec:crisis-deuda}); hoy pesan poco en los países pobres.
\item \emph{Los bonos}: desde el Plan \autor{Brady} (1989), los países emergentes y, desde los 2000, también algunos de renta baja (Ghana, Zambia o Kenia) se financian emitiendo bonos en los mercados internacionales.
  \begin{itemize}
  \item \emph{Ventajas}: amplían la base de inversores y alargan los plazos.
  \item \emph{Inconvenientes}: los tenedores son muchos y están dispersos, lo que complica reestructurar (\ref{sec:marco-deuda}), y los tipos son altos y suben de golpe cuando cambia el humor de los mercados.
  \item \emph{El pecado original}: la mayoría de estos países no puede endeudarse en el exterior en su propia moneda, de modo que una depreciación encarece la deuda\nocite{EichengreenHausmann1999,EHP2005PecadoOriginal} \vertema{3.B.17}.
  \end{itemize}
\item \emph{La financiación oficial}:
  \begin{itemize}
  \item \emph{bilateral}, de los gobiernos y de sus agencias de crédito a la exportación, con China como primer acreedor bilateral de los países pobres desde los 2010 (\ref{sec:crisis-deuda});
  \item \emph{multilateral}, del FMI (con condicionalidad y, para los países de renta baja, en condiciones de favor) y del Banco Mundial y los bancos regionales de desarrollo \vertema{3.B.22} \vertema{3.B.30}.
  \item \emph{Ventajas}: plazos largos, tipos bajos y, a menudo, anticíclica: llega cuando los mercados se cierran.
  \item \emph{Inconvenientes}: viene con condiciones de política económica y no basta para las necesidades de inversión.
  \end{itemize}
\end{itemize}

\subsubsection{Los recursos internos y la financiación innovadora}\label{sec:innovadora}

\begin{itemize}
\item \emph{Los recursos internos}: desde Addis Abeba (2015), la agenda pone en primer lugar los impuestos del propio país; la cooperación fiscal internacional, contra la erosión de bases imponibles, ayuda a recaudarlos \vertema{4.B.18}.
\item \emph{La financiación combinada} (\emph{blended finance}): usar fondos públicos o de la ayuda (donaciones, garantías, tramos que asumen las primeras pérdidas) para atraer inversión privada a proyectos que por sí solos no la atraerían. Es la idea de pasar «de miles de millones a billones» que propusieron en 2015 los bancos multilaterales y el FMI \vertema{3.B.30}.
\item \emph{Las asociaciones público-privadas}: el sector privado construye y gestiona la infraestructura y comparte los riesgos con el Estado, con el peligro de crear pasivos contingentes que no se ven en la deuda.
\item \emph{Los bonos temáticos y los canjes}:
  \begin{itemize}
  \item los bonos verdes, sociales y ligados a la sostenibilidad comprometen el destino de los fondos o los objetivos del emisor;
  \item en los \concepto{canjes de deuda} por desarrollo o por naturaleza, el acreedor perdona deuda a cambio de que el deudor gaste lo ahorrado en un fin acordado \vertema{4.A.24}.
  \end{itemize}
\item \emph{La financiación climática}: el objetivo de 100.000 millones de dólares anuales, fijado para 2020, solo se alcanzó en 2022, con 115.900 millones\nocite{OCDE2024Financiacion}; en Bakú (COP29, 2024) los países desarrollados se comprometieron a liderar la movilización de al menos 300.000 millones al año en 2035, dentro de un llamamiento a llegar a 1,3 billones de todas las fuentes\nocite{CMNUCC2024NCQG} \vertema{3.B.31}.
\end{itemize}

\begin{table}[H]
\centering
\caption{Las fuentes de la financiación exterior de un país en desarrollo}\label{tab:fuentes-desarrollo}
\small
\begin{tabular}{>{\raggedright\arraybackslash}m{2.9cm}>{\raggedright\arraybackslash}m{2.5cm}>{\raggedright\arraybackslash}m{3.4cm}>{\raggedright\arraybackslash}m{2.6cm}>{\raggedright\arraybackslash}m{2.6cm}}
\toprule
\rowcolor{tceeazulpalido} Fuente & ¿Genera deuda? & Estabilidad & Quién asume el riesgo & Además aporta \\
\midrule
Inversión directa & No (salvo préstamos entre empresas del grupo) & La más estable; sus rentas caen en las malas épocas & El inversor & Tecnología, gestión, mercados \\
Inversión de cartera en acciones & No & Volátil, pero sus rentas caen en las crisis & El inversor & Profundidad financiera \\
Remesas & No & Estables, a menudo anticíclicas & — & Renta de los hogares \\
Donaciones oficiales & No & Dependen del presupuesto del donante & El donante & Asistencia técnica \\
Préstamos bancarios y bonos & Sí & Volátil, con pagos fijos aunque el país esté en crisis; riesgo de parada súbita & El país (más aún en divisas) & Plazos largos (los bonos) \\
Préstamos oficiales & Sí & Estables, a menudo anticíclicos & El país, en condiciones de favor & Condicionalidad, asistencia técnica \\
\bottomrule
\end{tabular}
\fuente{Elaboración propia a partir de \textcite{FMI2014DeudaExterna} y \textcite{Calvo1998SuddenStops}.}
\end{table}

\begin{itemize}
\item \emph{Recapitulación}: la financiación exterior se justifica por la brecha de ahorro y de divisas y, mejor, por el enfoque intertemporal; pero el capital va poco a los pobres, y lo que más les llega son remesas e inversión directa, que no crean deuda. \uline{Cuanto más estable es la financiación y más riesgo asume quien la da, menos vulnerable es el país que la recibe.}
\item La deuda es la forma de financiación que más riesgo deja en el país que la recibe: pasamos a su problema.
\end{itemize}

% ==========================================================================
\section{El problema de la deuda externa}\label{sec:deuda}
\minutos{9}

\begin{itemize}
\item Una vez vistas las fuentes, nos centramos en las que crean deuda, porque son las que pueden convertir la financiación del desarrollo en su freno.
\item \emph{Conclusión anticipada}: la deuda es sostenible si el país crece más de lo que le cuesta; un Estado paga más por lo que perdería si no pagara que por obligación, y cuando la deuda es excesiva perjudica a la vez al deudor y al acreedor, por eso las crisis acaban en quitas.
\end{itemize}

% --------------------------------------------------------------------------
\subsection{Sostenibilidad, impago y sobreendeudamiento}\label{sec:sostenibilidad-impago}

\begin{itemize}
\item La deuda plantea 3 preguntas, que ordenan este subapartado: ¿puede pagarla el país?, ¿por qué la paga si nadie puede obligarle?, y ¿qué pasa cuando debe más de lo que podrá pagar?
\end{itemize}

\subsubsection{La sostenibilidad de la deuda externa}\label{sec:sostenibilidad}

\paragraph{a) Idea}

\begin{itemize}
\item \emph{La pregunta}: ¿cuándo deja de poder pagarse la deuda con la que un país financió su desarrollo?
\item \emph{Los 3 conceptos} que hay que distinguir\nocite{MilesiFerrettiRazin1996}:
  \begin{itemize}
  \item la \concepto{liquidez} es poder refinanciar lo que vence cada periodo: un país solvente puede tener un problema de liquidez si se le cierra el crédito;
  \item la \concepto{solvencia} es que el valor actual de los superávits futuros cubra la deuda exterior neta;
  \item la \concepto{sostenibilidad} es poder seguir pagando el servicio de la deuda (principal e intereses) sin financiación excepcional ni un ajuste brusco e irreal de ingresos o gastos\nocite{FMI2002Sustainability}.
  \end{itemize}
\item \emph{La crisis de deuda externa} llega cuando el país no puede atender el servicio de su deuda con el exterior \vertema{3.B.17}.
\end{itemize}

\paragraph{b) Modelo}

\subparagraph{i) Supuestos}

\begin{itemize}
\item \emph{Variables}: $b_t$ es la deuda exterior neta y $d_t$ el déficit comercial (el déficit exterior sin las rentas de la inversión), los 2 en porcentaje del PIB; $r^*$ es el tipo real que paga la deuda y $\gamma_Y$ el crecimiento real, los 2 constantes.
\item \emph{Moneda}: la deuda está en moneda nacional, o el tipo de cambio real no cambia.
\item \emph{Aviso de notación}: aquí $b$ es la deuda exterior y $d$ el déficit comercial; en la deuda pública, $d$ es la deuda y $g-\tau$ el déficit primario \vertema{3.A.39}.
\end{itemize}

\subparagraph{ii) Desarrollo}

\begin{itemize}
\item \emph{La dinámica}: la deuda de este año es la del año pasado, que crece con su tipo real y se diluye con el crecimiento, más el déficit comercial del año \vertema{3.B.12}:
  \funcionclave[eq:deuda-externa]{b_t=\frac{1+r^*}{1+\gamma_Y}\,b_{t-1}+d_t}
\item \emph{El saldo que estabiliza la deuda}: la ratio no cambia ($b_t=b_{t-1}=b$) si
  \funcionclave[eq:saldo-estabilizador]{-d=\frac{r^*-\gamma_Y}{1+\gamma_Y}\,b}
  \begin{itemize}
  \item \emph{En palabras}: \uline{el superávit comercial tiene que cubrir el coste de la deuda viva neto de lo que la diluye el crecimiento}.
  \item \emph{Un ejemplo}: con una deuda exterior del 60~\% del PIB, un tipo real del 5~\% y un crecimiento del 2~\%, hace falta un superávit comercial de casi el 1,8~\% del PIB ($0{,}03\cdot60/1{,}02$) solo para que la deuda no suba.
  \end{itemize}
\item \emph{Interpretación gráfica} (gráfico~\ref{fig:dinamica-deuda-externa}): en el plano $(b_{t-1},b_t)$, \eqref{eq:deuda-externa} es una recta de pendiente $(1+r^*)/(1+\gamma_Y)$ que corta al eje vertical en el déficit comercial. Con $r^*>\gamma_Y$ la pendiente es mayor que 1:
  \begin{itemize}
  \item con déficit comercial, la deuda crece sin límite: es la \concepto{bola de nieve};
  \item con superávit comercial, la recta corta a la bisectriz en $\overline b$, un umbral inestable: por debajo, la deuda baja; por encima, se dispara.
  \end{itemize}
\end{itemize}

\begin{figure}[H]
\centering
\caption{La dinámica de la deuda exterior con $r^*>\gamma_Y$}\label{fig:dinamica-deuda-externa}
\grafico[0.86\textwidth]{dinamica-deuda-externa}
\fuente{Elaboración propia a partir de \textcite[cap.~2]{ObstfeldRogoff1996Foundations}. A la izquierda, con déficit comercial, la deuda crece sin límite; a la derecha, con superávit comercial, $\overline b$ es un umbral inestable.}
\end{figure}

\subparagraph{iii) Implicaciones}

\begin{itemize}
\item \emph{La deuda concesional} de los países más pobres, a tipos muy bajos y plazos muy largos, suele tener $r^*<\gamma_Y$: la recta es más plana que la bisectriz y la ratio converge aunque haya un déficit comercial moderado. Por eso la ayuda en forma de préstamo apenas pone en riesgo la sostenibilidad y la deuda comercial sí.
\item \emph{Si la deuda está en divisas}, una depreciación real la encarece, porque $\gamma_q$ multiplica a $b_{t-1}$ (anexo~\ref{anexo:deuda}). Es el \concepto{pecado original}, la incapacidad de un país de endeudarse en el exterior en su propia moneda\nocite{EichengreenHausmann1999} \vertema{3.B.17}: la inflación propia no ayuda a pagar.
\item \emph{Los determinantes}: la sostenibilidad empeora con más deuda, menos crecimiento, tipos más altos, más deuda en divisas y a corto plazo y peor acceso a los mercados; mejora con una buena relación real de intercambio y con pasivos en forma de acciones o inversión directa.
\item \emph{Los indicadores} (tabla~\ref{tab:indicadores-deuda}): la deuda sobre el PIB no basta, porque no todo el PIB se puede vender fuera para pagar; en los países pobres pesan más los que miden la deuda frente a las exportaciones y a los ingresos fiscales.
\end{itemize}

\begin{table}[H]
\centering
\caption{Indicadores de la deuda externa de un país en desarrollo}\label{tab:indicadores-deuda}
\small
\begin{tabular}{>{\raggedright\arraybackslash}m{4.3cm}>{\raggedright\arraybackslash}m{5.0cm}>{\raggedright\arraybackslash}m{4.6cm}}
\toprule
\rowcolor{tceeazulpalido} Indicador & Qué mide & Referencia \\
\midrule
Deuda externa bruta / PIB & La carga bruta: todos los pasivos exigibles frente a no residentes & Sin umbral universal \\
Deuda externa / exportaciones & La deuda frente a la capacidad de generar divisas & PPME: 150~\% en valor actual \\
Servicio de la deuda / exportaciones & La liquidez: cuánto de lo que se ingresa en divisas se va en pagar & Renta baja: 10, 15 o 21~\% \\
Servicio de la deuda / ingresos fiscales & La carga para el presupuesto & Renta baja: 14, 18 o 23~\% \\
Tasa de renovación (desembolsos / amortizaciones) & Si los acreedores siguen prestando & Menor que 1: la financiación se retira \\
Composición por moneda, plazo y acreedor & Los descalces y lo difícil que será reestructurar & — \\
\bottomrule
\end{tabular}
\fuente{Elaboración propia a partir de \textcite{FMI2014DeudaExterna}, \textcite{FMI2018LICDSF} y \textcite{FMI1999HIPCReforzada}. Los umbrales de renta baja son los del marco del FMI y el Banco Mundial para una capacidad de endeudamiento débil, media o fuerte (anexo~\ref{anexo:umbrales}).}
\end{table}

\paragraph{c) Valoración}

\begin{itemize}
\item \emph{Sus límites}: no hay un umbral que valga para todos, porque los países con un historial de impagos entran en crisis con deudas mucho menores (la intolerancia a la deuda)\nocite{Reinhart2003Intolerancia}; y con equilibrios múltiples, el miedo de los acreedores sube la prima de riesgo y puede volver insostenible una deuda que no lo era \vertema{3.A.39}.
\item \emph{Lo que no dice}: mide si el país \emph{puede} pagar, pero un Estado puede decidir no hacerlo. ¿Por qué paga, entonces?
\end{itemize}

\subsubsection{¿Por qué pagan los Estados?}\label{sec:impago}

\paragraph{a) Idea}

\begin{itemize}
\item \emph{El problema}: una empresa que no paga quiebra y sus acreedores se quedan con sus activos; un Estado que no paga no quiebra, y casi nada suyo se puede embargar.
\item \emph{La consecuencia}: \uline{la deuda soberana depende de la voluntad de pagar, no solo de la capacidad}. Hace falta un modelo que explique por qué un país decide pagar.
\end{itemize}

\paragraph{b) Modelo: la reputación}

\subparagraph{i) Supuestos}

\begin{itemize}
\item \emph{El país}: tiene una renta que fluctúa y quiere suavizar su consumo, pidiendo prestado en los años malos y devolviendo en los buenos\nocite{EatonGersovitz1981}.
\item \emph{El castigo}: si no paga, los acreedores le cierran el crédito para siempre y queda en autarquía financiera.
\end{itemize}

\subparagraph{ii) Desarrollo}

\begin{itemize}
\item \emph{La decisión}: en cada periodo el país compara el valor de pagar y seguir en el mercado con el de no pagar y vivir sin crédito.
\item \emph{El resultado}: paga mientras lo que perdería por no poder suavizar su consumo en el futuro supere lo que se ahorra dejando de pagar hoy. Eso fija un \emph{techo de deuda} por encima del cual nadie le presta.
\end{itemize}

\subparagraph{iii) Implicaciones}

\begin{itemize}
\item \emph{Racionamiento del crédito}: los países pobres y volátiles se endeudan menos de lo que querrían, aunque tengan proyectos rentables. Es otro fallo del mercado de capitales detrás de la paradoja de \autor{Lucas} (\ref{sec:lucas}).
\item \emph{El momento del impago}: en las versiones cuantitativas del modelo, el impago llega en las recesiones, cuando pagar cuesta más, y los diferenciales suben antes\nocite{Arellano2008}.
\end{itemize}

\paragraph{c) Extensiones: las sanciones}

\begin{itemize}
\item \emph{La crítica}: si, tras el impago, el país puede seguir ahorrando fuera, suaviza su consumo sin pedir prestado, y la amenaza de cerrarle el crédito deja de bastar. \uline{La reputación sola no sostiene la deuda}\nocite{BulowRogoff1989AER}.
\item \emph{Lo que la sostiene}: las sanciones directas. Los acreedores pueden embargar los bienes y los pagos del país fuera de sus fronteras y entorpecer su comercio, así que las deudas se renegocian bajo esa amenaza, y lo que se cobra depende del poder de negociación de cada parte\nocite{BulowRogoff1989JPE}.
\end{itemize}

\paragraph{d) Evidencia empírica}

\begin{itemize}
\item \emph{Los costes del impago existen}: la producción cae en el año del impago y el acceso a los mercados se corta, aunque la penalización en los diferenciales dura poco\nocite{BorenszteinPanizza2009}.
\item \emph{Cuanto mayor la quita, mayor el castigo}: tras las reestructuraciones con quitas grandes, los diferenciales son más altos y la exclusión de los mercados más larga\nocite{CrucesTrebesch2013}.
\item \emph{El impago es recurrente}: muchos países han dejado de pagar varias veces en 2 siglos y han vuelto después al mercado\nocite{ReinhartRogoff2009This}.
\end{itemize}

\paragraph{e) Valoración}

\begin{itemize}
\item \emph{Lo que explican}: por qué se presta poco a los países pobres y por qué las crisis acaban en renegociaciones.
\item \emph{Lo que no}: suponen que una deuda mayor es siempre mejor para el acreedor. ¿Y si no lo es?
\end{itemize}

\subsubsection{El sobreendeudamiento}\label{sec:sobreendeudamiento}

\paragraph{a) Idea}

\begin{itemize}
\item \emph{El problema}: si un país debe más de lo que podrá pagar, cada euro que gane de más irá a los acreedores, y deja de invertir y de ajustar.
\item \emph{La consecuencia}: \uline{la deuda excesiva reduce lo que se puede cobrar}. Es el \concepto{sobreendeudamiento} (\emph{debt overhang}), que formalizaron \autor{Paul Krugman} (1988) y \autor{Jeffrey Sachs} (1989)\nocite{Krugman1988Overhang,Sachs1989Overhang}.\footnote{\autor{Paul Krugman} recibió el Premio Nobel de Economía en 2008 «por su análisis de los patrones del comercio y de la localización de la actividad económica»\nocite{Nobel2008Krugman}: lo premiado es su obra de comercio internacional, no la de la deuda.}
\item Con una deuda excesiva, el acreedor que exige cobrarla entera acaba cobrando menos: «quien presta, no cobra; y si cobra, no todo».
\end{itemize}

\paragraph{b) Modelo}

\subparagraph{i) Supuestos}

\begin{itemize}
\item \emph{El pago}: el país paga lo que debe o, si no le llega, la parte de su producción que los acreedores pueden obligarle a entregar.
\item \emph{La producción}: depende de la inversión y del esfuerzo de ajuste que haga hoy, y el futuro es incierto.
\item \emph{Notación}: $B$ es la deuda nominal (lo que se debe) y $V(B)$ el valor esperado de lo que cobran los acreedores.
\end{itemize}

\subparagraph{ii) Desarrollo}

\begin{itemize}
\item \emph{Con poca deuda}, el país la paga entera en todos los escenarios y se queda con todo lo que gane de más: invierte lo eficiente y $V(B)=B$.
\item \emph{Con más deuda que la que puede pagar en los malos escenarios}, parte del rendimiento de invertir hoy irá mañana a los acreedores. Funciona como un impuesto sobre la inversión: el país invierte y ajusta menos (anexo~\ref{anexo:deuda}).
\item \emph{La curva de Laffer de la deuda} relaciona el cobro esperado con la deuda nominal\nocite{Krugman1989MarketBased} (gráfico~\ref{fig:laffer-deuda}):
  \begin{itemize}
  \item hasta $\underline B$, la deuda se cobra entera y la curva sigue la bisectriz;
  \item en el tramo creciente, más deuda nominal es más cobro esperado, aunque cada euro valga menos: el precio de la deuda en el mercado secundario, $V(B)/B$, la pendiente del rayo desde el origen, baja;
  \item a la derecha del máximo, $B^{\max}$, más deuda es menos cobro: el desincentivo pesa más que el derecho a cobrar.
  \end{itemize}
\end{itemize}

\begin{figure}[H]
\centering
\caption{La curva de Laffer de la deuda}\label{fig:laffer-deuda}
\grafico[0.66\textwidth]{laffer-deuda}
\fuente{Elaboración propia a partir de \textcite{Krugman1988Overhang} y \textcite{Sachs1989Overhang}. Hasta $\underline B$ se cobra toda la deuda; a la derecha del máximo, una quita de $B_1$ a $B_2$ aumenta lo que cobran los acreedores. La pendiente del rayo es el precio de la deuda en el mercado secundario.}
\end{figure}

\subparagraph{iii) Implicaciones}

\begin{itemize}
\item \emph{La quita}: \uline{a la derecha del máximo, perdonar deuda beneficia a los 2}: el deudor vuelve a invertir y el acreedor cobra más. Es la justificación del Plan \autor{Brady} y de las condonaciones a los países pobres (\ref{sec:crisis-deuda}).
\item \emph{La acción colectiva}: cada acreedor prefiere que perdonen los demás y cobrar él entero (el acreedor que resiste, o \emph{holdout}), y eso obliga a coordinarlos (\ref{sec:marco-deuda}).
\end{itemize}

\paragraph{c) Evidencia empírica}

\begin{itemize}
\item \emph{La recompra de Bolivia} (1988): recompró casi la mitad de su deuda bancaria a unos 11 centavos por dólar; el valor de mercado de lo que quedaba apenas bajó, así que la ganancia fue sobre todo de los acreedores. Estaba donde la curva es casi plana: reducir la deuda apenas bajaba lo que se esperaba cobrar, así que la recompra al precio medio fue un regalo a los acreedores\nocite{BulowRogoff1988Buyback}.
\item \emph{Deuda y crecimiento}: en los países en desarrollo, a partir de cierto nivel, más deuda externa va unida a menos crecimiento\nocite{PattilloPoirsonRicci2002}.
\end{itemize}

\paragraph{d) Valoración}

\begin{itemize}
\item \emph{Sus límites}: es difícil saber en qué tramo está cada país, y perdonar crea riesgo moral.
\item \emph{Su legado}: cambió la política de la deuda, de prestar más a perdonar parte.
\end{itemize}

\begin{itemize}
\item \emph{Recapitulación}: la deuda es sostenible si el crecimiento cubre su coste; los Estados pagan sobre todo por miedo a las sanciones; y cuando la deuda es excesiva, perdonar parte puede convenir a todos. Con estas 3 ideas leemos ahora las crisis.
\end{itemize}

% --------------------------------------------------------------------------
\subsection{Las crisis de la deuda y sus respuestas}\label{sec:crisis-deuda}

\begin{itemize}
\item Ya tenemos la teoría; pasamos a la historia, en 3 episodios, cada uno con su respuesta.
\end{itemize}

\subsubsection{La crisis de la deuda de los 80}\label{sec:crisis-80}

\begin{itemize}
\item \emph{El origen}: en los 70, los bancos de los países ricos recibieron los depósitos de los exportadores de petróleo (los petrodólares) y los prestaron a América Latina en dólares y a tipo variable, con tipos reales bajos o negativos.
\item \emph{El detonante}: la subida de tipos de la Reserva Federal de \autor{Paul Volcker} (1979-1982) apreció el dólar, encareció la deuda en dólares de América Latina y, con la recesión mundial y la caída de las materias primas, desencadenó la crisis de la deuda de 1982 \vertema{3.B.17}. En agosto de 1982, México anunció que no podía pagar, y la crisis se extendió a casi toda la región.
\item \emph{El riesgo para los bancos}: en 1982, la deuda latinoamericana en manos de los 9 mayores bancos de Estados Unidos equivalía al 176~\% de su capital, de modo que un impago generalizado los habría dejado insolventes\nocite{SimsRomero2013LatAm}.
\item \emph{La primera respuesta}, la \concepto{financiación concertada} (1982-1985), trató la crisis como un problema de liquidez:
  \begin{itemize}
  \item el FMI prestó con programas de ajuste y el Club de París aplazó la deuda oficial;
  \item los bancos refinanciaron y pusieron «dinero nuevo» en proporción a lo que ya habían prestado;
  \item ganó tiempo para los bancos, pero no redujo la deuda.
  \end{itemize}
\item \emph{El Plan \autor{Baker}} (1985), que lleva el nombre del secretario del Tesoro de Estados Unidos, \autor{James Baker}, propuso «ajuste con crecimiento»: reformas estructurales a cambio de más préstamos de los bancos y de los organismos multilaterales. El dinero nuevo no llegó, y en 1987 Brasil suspendió pagos.
\item \emph{El Plan \autor{Brady}} (1989), que lleva el nombre de su sucesor, \autor{Nicholas Brady}, reconoció que era un problema de solvencia y que la deuda tenía que bajar:
  \begin{itemize}
  \item cambió préstamos bancarios por bonos a largo plazo, con quitas y con la garantía de bonos del Tesoro de Estados Unidos (los \concepto{bonos Brady}); los bancos elegían entre bonos con quita en el principal y bonos sin quita pero con un tipo bajo;
  \item lo usaron 17 países hasta 1998, empezando por México\nocite{FMI2023Brady};
  \item funcionó: los países que obtuvieron quitas en el principal crecieron más y redujeron después su deuda, al contrario que los que solo alargaron plazos\nocite{ReinhartTrebesch2016}. Es la curva de Laffer de la deuda en la práctica, y el origen del mercado de bonos de los países emergentes.
  \end{itemize}
\item \emph{El legado}: la «década perdida» de América Latina y una condicionalidad que se resumió en el \concepto{Consenso de Washington}\nocite{Williamson1990}: estabilidad macroeconómica, apertura y liberalización a cambio de financiación \vertema{3.B.28}.
\end{itemize}

\subsubsection{Los países más pobres: de los aplazamientos a la condonación}\label{sec:ppme}

\begin{itemize}
\item \emph{Las crisis de los 90 en los países de renta media} (México, Asia, Rusia, Brasil, Argentina) fueron sobre todo crisis de la cuenta financiera: paradas súbitas, descalces de monedas y crisis gemelas \vertema{3.B.17}. El Club de París respondió en 2003 con el \emph{enfoque de Evian}, un trato adaptado a cada país deudor, con quitas solo si la deuda es insostenible\nocite{ClubParis2026Miembros}.
\item \emph{La deuda de los países más pobres}, sobre todo en África subsahariana, era oficial (de gobiernos y organismos multilaterales), y los aplazamientos del Club de París solo la acumulaban. Tras la campaña del Jubileo 2000, se pasó a perdonarla.
\item \emph{La Iniciativa para los Países Pobres Muy Endeudados} (PPME; en inglés, HIPC), del FMI y el Banco Mundial (1996, reforzada en 1999), reduce la deuda hasta niveles sostenibles a cambio de reformas y de una estrategia contra la pobreza, en 2 etapas\nocite{FMI2023HIPC,FMI1999HIPCReforzada}:
  \begin{enumerate}
  \item en el \emph{punto de decisión}, el país muestra que cumple los requisitos (entre ellos, una deuda superior al 150~\% de sus exportaciones en valor actual) y empieza a recibir un alivio provisional;
  \item en el \emph{punto de culminación}, tras aplicar las reformas, recibe la reducción definitiva, también de los organismos multilaterales, que por primera vez renuncian a cobrar lo que se les debe.
  \end{enumerate}
\item \emph{La Iniciativa para el Alivio de la Deuda Multilateral} (IADM; en inglés, MDRI), de 2005, condona el 100~\% de la deuda con el FMI, la Asociación Internacional de Fomento (AIF) del Banco Mundial y el Fondo Africano de Desarrollo a los países que llegan al punto de culminación\nocite{FMI2023HIPC} \vertema{3.B.30}.
\item \emph{El resultado}: de los 39 países de la PPME, 37 han llegado al punto de culminación (el último, Somalia, en diciembre de 2023, cuya deuda externa baja del 64~\% del PIB en 2018 a menos del 6~\%), Sudán está en el de decisión desde 2021 y Eritrea no ha llegado a él\nocite{FMI2023HIPC,FMI2023Somalia}. Los requisitos están en el anexo~\ref{anexo:umbrales}.
\item \emph{La valoración}: redujeron mucho la deuda y liberaron gasto social, pero muchos países volvieron a endeudarse en pocos años, ahora con acreedores nuevos: los bonistas y China.
\end{itemize}

\subsubsection{De la pandemia a hoy}\label{sec:deuda-hoy}

\begin{itemize}
\item \emph{La nueva ola de deuda}: con los tipos bajos de los 2010, los países de renta baja emitieron bonos y pidieron mucho a China, que se convirtió en el mayor acreedor oficial de los países en desarrollo, con préstamos poco transparentes y a menudo garantizados con materias primas\nocite{HornReinhartTrebesch2021}.
\item \emph{La pandemia}: la \concepto{Iniciativa de Suspensión del Servicio de la Deuda} (ISSD) del G-20 aplazó más de 12.000 millones de dólares de pagos de casi 50 países entre mayo de 2020 y diciembre de 2021, pero solo de la deuda oficial bilateral\nocite{G202020DSSI,BancoMundial2022DSSI} \vertema{3.B.19}.
\item \emph{El Marco Común} para el tratamiento de la deuda más allá de la ISSD (noviembre de 2020) sienta por primera vez en la misma mesa al Club de París y a los acreedores del G-20 que no son miembros, como China, la India o Arabia Saudí\nocite{G202020MarcoComun}:
  \begin{itemize}
  \item el tratamiento parte de un programa del FMI y de su análisis de sostenibilidad, y los acreedores privados tienen que dar un trato comparable;
  \item lo han pedido 4 países: Chad, el primero en cerrar un acuerdo con sus acreedores oficiales (2022), Zambia, Ghana y Etiopía\nocite{Francia2026MarcoComun};
  \item ha sido lento, con procesos de años, y no llega a los países de renta media, que reestructuran fuera de él.
  \end{itemize}
\item \emph{La Mesa Redonda Global sobre la Deuda Soberana} (desde febrero de 2023), del FMI, el Banco Mundial y la presidencia del G-20, reúne a deudores y acreedores para acordar reglas comunes y acortar los plazos\nocite{FMI2023GSDR}.
\item \emph{Hoy}:
  \begin{itemize}
  \item la deuda externa de los países de renta baja y media llegó en 2025 a un récord de 9,4 billones de dólares, un 5,5~\% más que un año antes, y pagaron 1,4 billones de principal e intereses, el 13,1~\% de sus exportaciones\nocite{BancoMundial2026IDR,BancoMundial2026IDS};
  \item entre 2022 y 2025 pagaron 436.000 millones de dólares más de lo que recibieron en financiación nueva, aunque en 2025 los flujos netos de deuda a largo plazo volvieron a ser positivos\nocite{BancoMundial2026IDR};
  \item a 31 de marzo de 2026, de los países de renta baja que analiza el marco del FMI y el Banco Mundial, 9 estaban en sobreendeudamiento y 23 en riesgo alto\nocite{FMI2026ListaDSA}.
  \end{itemize}
\item \emph{Recapitulación}: cada crisis cambió la respuesta, de prestar más (1982) a perdonar (Brady, PPME) y a coordinar a acreedores cada vez más distintos (Marco Común). Eso pide un marco institucional, que vemos a continuación.
\end{itemize}

\begin{anotaciones}
\begin{itemize}
\item \emph{Los impagos recientes}: Zambia (noviembre de 2020), Sri Lanka (2022) y Ghana (2022) dejaron de pagar sus bonos; Zambia y Ghana reestructuraron dentro del Marco Común; Etiopía sigue en ello, y Sri Lanka, de renta media, lo hizo fuera de él\nocite{Francia2026MarcoComun,FMI2026ListaDSA}.
\item \emph{Los casos del Marco Común} se siguen en la página del Club de París\nocite{Francia2026MarcoComun} y en los informes de la Mesa Redonda Global.
\item \emph{La cronología completa} de la deuda y de la agenda de la financiación del desarrollo está en el anexo~\ref{anexo:cronologia}.
\end{itemize}
\end{anotaciones}

% --------------------------------------------------------------------------
\subsection{El marco institucional y las soluciones}\label{sec:marco-deuda}

\begin{itemize}
\item \emph{El punto de partida}: no hay un tribunal de quiebras para los Estados. La deuda soberana se resuelve en foros informales, con el FMI como árbitro técnico, y el marco se ordena por quién es el acreedor.
\end{itemize}

\subsubsection{El FMI y el análisis de sostenibilidad}\label{sec:dsa}

\begin{itemize}
\item \emph{El FMI} presta con condiciones a los países con problemas de balanza de pagos, con servicios en condiciones de favor para los de renta baja, y no presta si la deuda no es sostenible: su programa es la puerta de entrada de cualquier reestructuración \vertema{3.B.22}.
\item \emph{Su análisis de sostenibilidad} tiene 2 marcos \vertema{3.B.12}:
  \begin{itemize}
  \item el de los países con acceso a los mercados, centrado en la deuda pública;
  \item el \concepto{marco de sostenibilidad de la deuda para los países de renta baja} (\emph{low-income country debt sustainability framework}, LIC DSF), común con el Banco Mundial.
  \end{itemize}
\item \emph{El LIC DSF} clasifica el riesgo de sobreendeudamiento externo en bajo, moderado, alto o en sobreendeudamiento, comparando 4 indicadores con umbrales que dependen de la \emph{capacidad de endeudamiento} del país (débil, media o fuerte), como el valor actual de la deuda externa pública sobre el PIB (30, 40 o 55~\%) y sobre las exportaciones (140, 180 o 240~\%)\nocite{FMI2018LICDSF} (anexo~\ref{anexo:umbrales}).
  \begin{itemize}
  \item El resultado decide cuánto da el Banco Mundial en donaciones y cuánto en préstamos, y si el FMI puede aprobar un programa.
  \item Se revisó en septiembre de 2026, y el marco nuevo se aplicará desde el verano de 2027\nocite{FMI2026LICDSF}.
  \end{itemize}
\end{itemize}

\subsubsection{Los foros de acreedores}\label{sec:foros}

\begin{itemize}
\item \emph{El Club de París} reúne a los acreedores oficiales bilaterales desde su primera reunión con un deudor, Argentina, en 1956. Tiene 22 miembros permanentes, entre ellos España; los últimos en entrar fueron Corea y Brasil, en 2016\nocite{ClubParis2026Miembros,Francia2026ClubParis} \vertema{3.B.19}.
  \begin{itemize}
  \item \emph{Sus 6 principios}: solidaridad entre acreedores, consenso, intercambio de información, decisión caso por caso, condicionalidad (el deudor necesita un programa con el FMI) y comparabilidad de trato (no puede dar mejores condiciones a otros acreedores).
  \item \emph{La fecha de corte}: la deuda contraída después no se reestructura, para que el país pueda seguir recibiendo crédito nuevo.
  \end{itemize}
\item \emph{El Club de Londres} reúne a los bancos acreedores en un comité asesor para cada negociación, sin sede ni reglas fijas.
\item \emph{Los bonistas} son miles de tenedores dispersos que se organizan en comités \emph{ad hoc}, con el riesgo de los fondos que no aceptan la reestructuración y litigan para cobrar entero, como los que pleitearon contra Argentina tras su impago de 2001.
\end{itemize}

\subsubsection{Las reformas}\label{sec:reformas-deuda}

\begin{itemize}
\item \emph{La vía legal}: \autor{Anne Krueger}, primera subdirectora gerente del FMI, propuso en 2001-2002 un \concepto{mecanismo de reestructuración de la deuda soberana} (\emph{sovereign debt restructuring mechanism}, SDRM), una especie de quiebra para los Estados, con mayorías que vincularan a todos los acreedores y suspensión de los litigios\nocite{Krueger2002SDRM}. Estados Unidos y los mercados se opusieron, y se abandonó en 2003.
\item \emph{La vía contractual}: las \concepto{cláusulas de acción colectiva} (CAC) permiten que una mayoría cualificada de bonistas apruebe una reestructuración que vincula a todos. México las incluyó en 2003 y hoy son la norma en las emisiones internacionales \vertema{3.B.22}.
\item \emph{Las cláusulas de deuda resiliente al clima} suspenden los pagos tras una catástrofe; en Sevilla (2025) se lanzó una alianza de países y bancos de desarrollo para extenderlas\nocite{ONU2025FfD4Cierre} \vertema{3.B.31}.
\item \emph{Los canjes de deuda} por desarrollo o por naturaleza convierten pagos en gasto acordado en el propio país; España tiene programas de conversión de deuda \vertema{4.A.24}.
\end{itemize}

\begin{itemize}
\item \emph{Recapitulación}: la deuda externa puede financiar el desarrollo si crece menos de lo que crece el país; cuando no, el impago es una decisión y el exceso de deuda perjudica a todos, así que la solución ha acabado siendo la quita, cada vez más coordinada en un marco sin tribunal. \uline{La deuda que sirve al desarrollo es la que se toma a largo plazo, en moneda propia, y se reestructura a tiempo.}
\item Cuando la deuda no se puede pagar, la financiación que queda es la que no se devuelve: pasamos a la ayuda.
\end{itemize}

% ==========================================================================
\section{La ayuda al desarrollo}\label{sec:ayuda}
\minutos{8,5}

\begin{itemize}
\item Una vez vista la financiación que se devuelve, pasamos a la que se da en condiciones de favor: la ayuda.
\item \emph{Conclusión anticipada}: la ayuda es hoy una fuente menor y en retroceso, pero decisiva para los países más pobres; funciona mejor en programas concretos y evaluados que como relleno de una brecha de ahorro.
\end{itemize}

% --------------------------------------------------------------------------
\subsection{Concepto y medida}\label{sec:aod-concepto}

\begin{itemize}
\item \emph{De lo amplio a lo concreto}:
  \begin{itemize}
  \item la \concepto{cooperación al desarrollo} es todo lo que hacen los países, las instituciones y las organizaciones a favor del desarrollo de otros, incluidas las preferencias comerciales;
  \item la \concepto{ayuda al desarrollo} es la transferencia de recursos en condiciones más favorables que las del mercado, pública o privada (las organizaciones no gubernamentales y las fundaciones);
  \item la \concepto{ayuda oficial al desarrollo} (AOD) es la que cumple la definición del Comité de Ayuda al Desarrollo (CAD) de la OCDE \vertema{3.B.19}.
  \end{itemize}
\item \emph{La definición del CAD}: son AOD los flujos que cumplen 4 requisitos\nocite{OCDE2026DefinicionAOD}:
  \begin{enumerate}
  \item los dan organismos públicos;
  \item van a los países de la lista de receptores del CAD (los de renta baja y media, revisada cada 3 años) o a organismos multilaterales de desarrollo;
  \item tienen como objetivo principal el desarrollo económico y el bienestar del receptor (se excluye la ayuda militar);
  \item se dan en condiciones de favor (son \emph{concesionales}).
  \end{enumerate}
\item \emph{La donación equivalente}: desde los datos de 2018, un préstamo cuenta como AOD solo por su \concepto{elemento de donación}, la parte que es regalo, y solo si supera un mínimo: el 45~\% para los países menos adelantados y otros de renta baja, el 15~\% para los de renta media-baja y el 10~\% para los de renta media-alta y los organismos multilaterales\nocite{OCDE2026DefinicionAOD}. Antes se contaba el préstamo entero y se restaban sus devoluciones.
\item \emph{Los tipos de AOD}:
  \begin{itemize}
  \item \emph{por el canal}: \emph{bilateral}, de donante a receptor, y \emph{multilateral}, a través de los organismos internacionales;
  \item \emph{por el instrumento}: donaciones, préstamos en condiciones de favor, cooperación técnica y alivio de la deuda;
  \item \emph{por lo que llega al receptor}: una parte se gasta en el propio país donante (la acogida de refugiados en su primer año, los costes administrativos, las becas) y no sale de él.
  \end{itemize}
\end{itemize}

\begin{anotaciones}
\begin{itemize}
\item \emph{Cómo se calcula el elemento de donación}: es la diferencia entre el valor nominal del préstamo y el valor actual de lo que se devuelve (principal e intereses), en porcentaje del valor nominal. Cuanto más alto el tipo de descuento, más vale como regalo un préstamo barato; por eso el CAD descuenta más los préstamos a los países más pobres (tabla~\ref{tab:elemento-donacion}).
\item \emph{Antes de 2018}, el umbral era un 25~\% con un descuento del 10~\% para todos los receptores.
\end{itemize}
\begin{tablacaja}
\centering
\caption{Requisitos de concesionalidad de los préstamos para contar como AOD (desde 2018)}\label{tab:elemento-donacion}
\small
\begin{tabular}{>{\raggedright\arraybackslash}m{6.4cm}>{\centering\arraybackslash}m{3.0cm}>{\centering\arraybackslash}m{3.0cm}}
\toprule
\rowcolor{tceeazulpalido} Receptor & Elemento de donación mínimo & Tipo de descuento \\
\midrule
Países menos adelantados y otros de renta baja & 45~\% & 9~\% \\
Renta media-baja & 15~\% & 7~\% \\
Renta media-alta & 10~\% & 6~\% \\
Organismos multilaterales & 10~\% & 5 o 6~\% \\
\bottomrule
\end{tabular}
\fuente{Elaboración propia a partir de \textcite{OCDE2026DefinicionAOD}.}
\end{tablacaja}
\end{anotaciones}

\begin{itemize}
\item \emph{La ayuda ligada} es la que obliga al receptor a comprar en el país donante:
  \begin{itemize}
  \item \emph{sus costes}: encarece los proyectos entre un 15 y un 30~\%\nocite{Jepma1991Tying}, distorsiona su elección y funciona como una subvención a la exportación;
  \item \emph{su regulación}: el Acuerdo sobre Créditos a la Exportación con Apoyo Oficial de la OCDE (el \concepto{Consenso OCDE}), de 1978 y nacido de un consenso de 1976 entre menos países, pone límites a la ayuda ligada, que las reglas de Helsinki (1991) prohibieron para los proyectos comercialmente viables\nocite{OCDE1978Consenso} \vertema{4.A.24}; la Recomendación del CAD de 2001 desligó la ayuda a los países menos adelantados, y después a los PPME y a otros de renta baja\nocite{OCDE2018AyudaDesligada};
  \item \emph{por qué cuesta desligar}: es un dilema del prisionero entre donantes, porque el que desliga primero pierde contratos para sus empresas \vertema{3.A.14};
  \item \emph{hoy}: en 2022, el 81,5~\% de los compromisos bilaterales de AOD de los países del CAD estaba desligado\nocite{OCDE2026DAC7b}.
  \end{itemize}
\item \emph{Más allá de la AOD}: el \concepto{apoyo oficial total al desarrollo sostenible} (\emph{total official support for sustainable development}, TOSSD) mide también los flujos oficiales no concesionales, la cooperación Sur-Sur y la financiación privada que moviliza la intervención pública; es una de las fuentes del indicador 17.3.1 de los Objetivos de Desarrollo Sostenible\nocite{TOSSD2022}.
\end{itemize}

% --------------------------------------------------------------------------
\subsection{Evolución y agenda}\label{sec:aod-agenda}

\subsubsection{Del 0,7~\% a la eficacia de la ayuda}\label{sec:aod-eficacia-agenda}

\begin{itemize}
\item \emph{El origen}: tras el Plan Marshall y con la descolonización y la Guerra Fría, los países ricos crearon en 1961 el CAD para coordinar su ayuda.
\item \emph{El 0,7~\%}: la comisión del Banco Mundial que presidió \autor{Lester Pearson} propuso en 1969 que la AOD llegara al 0,7~\% de la renta de los donantes, y la Asamblea General de las Naciones Unidas lo adoptó en 1970 (resolución 2626)\nocite{OCDE2026Meta07}. Hoy se mide sobre la renta nacional bruta (RNB), y muy pocos donantes lo cumplen (gráfico~\ref{fig:aod-rnb}).
\item \emph{La fatiga de la ayuda} de los 90: tras la Guerra Fría la ayuda cayó hasta un mínimo del 0,21~\% de la RNB en 2001\nocite{OCDE2026DAC1}, y sus resultados decepcionaban.
\item \emph{La agenda de la eficacia}: la Declaración de París (2005) fijó 5 principios\nocite{OCDE2005Paris}:
  \begin{enumerate}
  \item \emph{apropiación}: el receptor decide su estrategia;
  \item \emph{alineamiento}: los donantes usan los sistemas del receptor;
  \item \emph{armonización}: los donantes se coordinan entre sí;
  \item \emph{gestión por resultados};
  \item \emph{mutua responsabilidad}.
  \end{enumerate}
\item \emph{Después}: Accra (2008) y Busan (2011), que abrió la agenda a los donantes emergentes y al sector privado.
\end{itemize}

\subsubsection{De la ayuda a la financiación del desarrollo}\label{sec:ffd}

\begin{itemize}
\item \emph{El cambio de enfoque}: las Naciones Unidas pasan de hablar de ayuda a hablar de todas las fuentes de la financiación del desarrollo, en 4 conferencias:
  \begin{enumerate}
  \item \emph{Monterrey} (2002): la responsabilidad es compartida; el país pone las políticas y los recursos internos, y los donantes, más ayuda, comercio y alivio de la deuda;
  \item \emph{Doha} (2008): lo reafirma en plena crisis financiera;
  \item \emph{Addis Abeba} (2015): el marco de financiación de la Agenda 2030, con los impuestos propios y la inversión privada en primer plano;
  \item \emph{Sevilla} (2025): el \emph{Compromiso de Sevilla}, aprobado sin Estados Unidos, que se había retirado de la negociación, y hecho suyo por la Asamblea General en la resolución 79/323; pide, entre otras cosas, triplicar la capacidad de préstamo de los bancos multilaterales y reformar la reestructuración de la deuda\nocite{ONU2025Sevilla,ONU2025SevillaNoticia}.
  \end{enumerate}
\end{itemize}

\subsubsection{Los datos y los recortes}\label{sec:aod-datos}

\begin{itemize}
\item \emph{La caída}: tras un máximo en 2023, la AOD bajó un 6~\% real en 2024\nocite{OCDE2025AODFinal2024} y un 23,1~\% en 2025, la mayor caída registrada, hasta 174.300 millones de dólares y el 0,26~\% de la renta nacional bruta de los donantes, el nivel más bajo desde 2004\nocite{OCDE2026AODPreliminar,OCDE2026DAC1}.
  \begin{itemize}
  \item Tres cuartas partes de la caída de 2025 son de Estados Unidos, que recortó su ayuda un 56,9~\%; Alemania pasó a ser, por primera vez, el mayor donante\nocite{OCDE2026AODPreliminar,OCDE2026DAC1}.
  \item Solo Noruega, Luxemburgo, Suecia y Dinamarca superan el 0,7~\%\nocite{OCDE2026AODPreliminar}.
  \item El Reino Unido bajará la suya del 0,5~\% al 0,3~\% de la renta nacional bruta en 2027 para pagar más defensa\nocite{ReinoUnido2025AOD}.
  \end{itemize}
\item \emph{España}: dio 4.554 millones de euros en 2025, el 0,27~\% de su renta nacional bruta y la cifra más alta en 15 años\nocite{OCDE2026DAC1,Exteriores2026AOD}, lejos del 0,7~\% que la Ley 1/2023 fija para 2030\nocite{BOE2023LeyCooperacion} \vertema{4.A.26}.
\item \emph{La Unión Europea}, con sus Estados miembros, es el mayor donante del mundo \vertema{3.B.42}.
\end{itemize}

\begin{figure}[H]
\centering
\caption{La AOD de los países del CAD y de España en porcentaje de la RNB, 1960-2025}\label{fig:aod-rnb}
\grafico[0.9\textwidth]{aod-rnb}
\fuente{Elaboración propia a partir de \textcite{OCDE2026DAC1}. Hasta 2017, desembolsos netos; desde 2018, donación equivalente (por eso la serie se corta). 2025, datos preliminares. El repunte de 2022 refleja el gasto en refugiados y la ayuda a Ucrania.}
\end{figure}

% --------------------------------------------------------------------------
\subsection{¿Funciona la ayuda?}\label{sec:aod-eficacia}

\begin{itemize}
\item El debate sobre la ayuda es, en el fondo, si da peces, si enseña a pescar o si quita las ganas de pescar.
\end{itemize}

\subsubsection{Por qué se da}\label{sec:aod-motivos}

\begin{itemize}
\item \emph{Los motivos de los donantes} no son solo el desarrollo: el pasado colonial y la afinidad política (cómo vota el receptor en las Naciones Unidas) explican el reparto de la ayuda tanto como la pobreza o las buenas políticas del receptor\nocite{AlesinaDollar2000}.
\item \emph{La justificación clásica}: cubrir la brecha de financiación (\ref{sec:dos-brechas}) o dar el \concepto{gran empujón} que saque a un país de una trampa de pobreza, coordinando inversiones que por separado no son rentables, como proponía \autor{Paul Rosenstein-Rodan} (1943) \vertema{3.B.28}. \textcite{Sachs2005} la actualiza: los países más pobres no salen solos de la trampa, y la ayuda tiene que ser mucho mayor.
\item \emph{La primera crítica}: \textcite{Bauer1972} sostenía que la ayuda refuerza a los gobiernos frente a la sociedad y que, si un proyecto es rentable, el mercado lo financia.
\end{itemize}

\subsubsection{La evidencia: ¿hace crecer la ayuda?}\label{sec:aod-evidencia}

\begin{itemize}
\item \emph{La paradoja micro-macro}: muchos proyectos concretos funcionan (vacunas, escuelas, carreteras), pero cuesta encontrar en los datos agregados que la ayuda haga crecer a los países.
\item \emph{Por qué podría no funcionar}:
  \begin{enumerate}
  \item la \emph{fungibilidad}: si el donante financia escuelas, el gobierno puede gastar en otra cosa lo que habría dedicado a ellas, de modo que la ayuda financia el gasto marginal, no el proyecto;
  \item la \emph{enfermedad holandesa de la ayuda}: la entrada de divisas aprecia el tipo de cambio real y resta competitividad a las manufacturas que se exportan\nocite{RajanSubramanian2011};
  \item las \emph{instituciones}: un gobierno que vive de la ayuda rinde cuentas a los donantes y no a sus ciudadanos\nocite{Deaton2013}.
  \end{enumerate}
\item \emph{El debate econométrico}:
  \begin{itemize}
  \item \textcite{BurnsideDollar2000}: la ayuda hace crecer a los países con buenas políticas fiscales, monetarias y comerciales, y apenas a los demás. Dio la base a la \emph{selectividad}: dar más a quien hace bien las cosas, como en el reparto por resultados de la AIF \vertema{3.B.30}.
  \item \textcite{EasterlyLevineRoodman2004}: con más países y más años, el resultado desaparece.
  \item \textcite{RajanSubramanian2008}: corrigiendo la causalidad inversa (se da más ayuda a quien va peor), no hay una relación robusta, ni positiva ni negativa, entre la ayuda y el crecimiento.
  \item \textcite{ClemensEtAl2012}: si se mira la ayuda que puede hacer crecer pronto (infraestructuras, apoyo a la producción) y se le da tiempo, los aumentos de la ayuda van seguidos de aumentos modestos de la inversión y del crecimiento.
  \item \textcite{GalianiEtAl2017}, con un cuasiexperimento: cuando un país supera el umbral de renta de la AIF, la ayuda que recibe cae; con ese salto, un punto más de ayuda sobre la renta nacional eleva el crecimiento por habitante en unos 0,35 puntos al año.
  \end{itemize}
\item \emph{La lectura}: \uline{la ayuda tiene un efecto positivo pero pequeño sobre el crecimiento, mucho menor que el que prometía la brecha de financiación}.
\end{itemize}

\subsubsection{El debate de fondo y la evaluación}\label{sec:aod-debate}

\begin{itemize}
\item \emph{Más ayuda}: para \autor{Jeffrey Sachs}, la ayuda ha fallado por escasa; con un gran empujón en salud, agricultura e infraestructuras, los más pobres pueden salir de la trampa\nocite{Sachs2005}.
\item \emph{Menos planes}: para \autor{William Easterly}, fracasan los \emph{planificadores} que diseñan desde arriba y funcionan los \emph{buscadores}, que prueban soluciones concretas, escuchan al terreno y responden ante los beneficiarios\nocite{Easterly2006}.
\item \emph{Menos ayuda}: para \autor{Angus Deaton}\footnote{\autor{Angus Deaton} recibió el Premio Nobel de Economía en 2015 «por su análisis del consumo, la pobreza y el bienestar»\nocite{Nobel2015Deaton}.}, la ayuda grande debilita el contrato entre el Estado y sus ciudadanos, y es mejor dejar de perjudicar desde fuera (patentes, comercio, armas)\nocite{Deaton2013}; para \autor{Dambisa Moyo}, hay que sustituirla por mercados y comercio\nocite{Moyo2009}.
\item \emph{La vía experimental}: \autor{Abhijit Banerjee}, \autor{Esther Duflo} y \autor{Michael Kremer}\footnote{\autor{Abhijit Banerjee}, \autor{Esther Duflo} y \autor{Michael Kremer} recibieron el Premio Nobel de Economía en 2019 «por su enfoque experimental para aliviar la pobreza global»\nocite{NobelPrize2019}.} cambian la pregunta de si la ayuda funciona por la de qué programa concreto funciona, y la responden con experimentos aleatorios\nocite{BanerjeeDuflo2011} \vertema{3.B.28}.
  \begin{itemize}
  \item Los programas evaluados (desparasitación, transferencias, programas de graduación) orientan hoy una parte de la ayuda.
  \item \emph{Su límite}: lo que funciona en un experimento no siempre vale en otro país ni a otra escala\nocite{Deaton2010} \vertema{4.B.6}.
  \end{itemize}
\end{itemize}

\begin{itemize}
\item \emph{Recapitulación}: la ayuda nació para cubrir una brecha que no era la importante y hoy cae cuando los más pobres pagan más por su deuda; se justifica por programas evaluados, por los bienes públicos mundiales (salud, clima) y por los países más pobres, para los que sigue siendo clave. \uline{Importa menos cuánta ayuda se da que cómo y a quién.}
\end{itemize}

% ==========================================================================
\section*{Conclusión}
\minutos{1,5}

\begin{itemize}
\item \textbf{\emph{Recapitulación}}:
  \begin{itemize}
  \item La teoría justifica que el capital vaya a los países pobres, pero va poco, y lo que más les llega son remesas e inversión directa, que no crean deuda.
  \item La deuda es sostenible si el crecimiento cubre su coste; los Estados pagan por lo que perderían si no, y si la deuda es excesiva, perdonar parte conviene a deudor y acreedor: así se cerraron la crisis de los 80 (Plan \autor{Brady}) y la de los países más pobres (PPME e IADM), y así se intenta cerrar la de hoy (Marco Común).
  \item La ayuda tiene un efecto pequeño sobre el crecimiento agregado, pero funciona en programas concretos y evaluados.
  \end{itemize}
\item \textbf{\emph{Relevancia}}:
  \begin{itemize}
  \item Con la ayuda en mínimos desde 2004 y la deuda externa en máximos, financiar los Objetivos de Desarrollo Sostenible depende de movilizar los impuestos propios y la inversión privada y de reformar la arquitectura de la deuda: es la agenda de Sevilla.
  \end{itemize}
\item \textbf{\emph{Extensiones}}:
  \begin{itemize}
  \item El enfoque intertemporal y la sostenibilidad exterior \vertema{3.B.12} y las crisis de balanza de pagos \vertema{3.B.17}.
  \item El FMI y sus programas \vertema{3.B.22} y el Banco Mundial y los bancos regionales de desarrollo \vertema{3.B.30}.
  \item Las teorías del desarrollo y la evaluación experimental \vertema{3.B.28} y la financiación climática \vertema{3.B.31}.
  \item La cooperación de la Unión Europea \vertema{3.B.42} y la de España, con sus instrumentos financieros \vertema{4.A.24} \vertema{4.A.26}.
  \end{itemize}
\item \textbf{\emph{Opinión}}:
  \begin{itemize}
  \item La lección de 70 años es que importa más \emph{cómo} se financia el desarrollo que \emph{cuánto}: financiación que comparte el riesgo, deuda a largo plazo y en moneda propia, reestructuraciones que lleguen a tiempo y ayuda evaluada.
  \end{itemize}
\item \textbf{\emph{Idea final}}:
  \begin{itemize}
  \item Volvamos al pez y a la caña. El capital no llega a los países pobres porque rinde menos de lo que parece y prestarles es arriesgado; llega caro cuando llega como deuda que se come lo que se pesca; y la ayuda que se da como pez no enseña a pescar.
  \item Lo que decide es lo que cada país hace con lo que recibe: instituciones, recursos propios y una deuda que se pueda pagar. Por eso la agenda de Monterrey a Sevilla pone primero los recursos internos: la mejor financiación del desarrollo es la que enseña a pescar.
  \end{itemize}
\end{itemize}

\clearpage
\appendix

% ==========================================================================
\section{Anexo: cronología}\label{anexo:cronologia}

\begin{table}[H]
\centering
\caption{Cronología de la deuda externa y de la financiación del desarrollo}\label{tab:cronologia-deuda}
\small
\begin{tabular}{>{\centering\arraybackslash}m{1.8cm}>{\raggedright\arraybackslash}m{12.2cm}}
\toprule
\rowcolor{tceeazulpalido} Año & Hito \\
\midrule
1956 & Primera reunión del Club de París (con Argentina) \\
1961 & Comité de Ayuda al Desarrollo de la OCDE \\
1966 & Las 2 brechas de \autor{Hollis Chenery} y \autor{Alan Strout} \\
1969-1970 & Informe \autor{Pearson}; la Asamblea General de las Naciones Unidas adopta el objetivo del 0,7~\% \\
1978 & Consenso OCDE sobre créditos a la exportación con apoyo oficial \\
1979-1982 & La Reserva Federal de \autor{Paul Volcker} sube los tipos \\
1982 & México no puede pagar: crisis de la deuda \\
1985 & Plan \autor{Baker} \\
1987 & Brasil suspende pagos \\
1988 & Recompra de deuda de Bolivia \\
1989 & Plan \autor{Brady} \\
1991 & Reglas de Helsinki sobre la ayuda ligada \\
1996 & Iniciativa para los PPME (reforzada en 1999) \\
2001 & Recomendación del CAD de desligar la ayuda a los países menos adelantados; impago de Argentina \\
2001-2003 & Propuesta del SDRM de \autor{Anne Krueger}, abandonada en 2003 \\
2002 & Conferencia de Monterrey \\
2003 & Enfoque de Evian del Club de París; cláusulas de acción colectiva de México \\
2005 & Declaración de París; IADM \\
2008 & Accra; Conferencia de Doha \\
2011 & Busan \\
2015 & Addis Abeba y Agenda 2030 \\
2018 & La AOD pasa a medirse en donación equivalente \\
2020 & ISSD (mayo) y Marco Común (noviembre); impago de Zambia \\
2022 & Impagos de Sri Lanka y Ghana; Chad, primer acuerdo del Marco Común \\
2023 & Mesa Redonda Global sobre la Deuda Soberana; Somalia, punto de culminación de la PPME \\
2025 & Compromiso de Sevilla; la AOD cae un 23,1~\% \\
2026 & Revisión del LIC DSF (se aplicará desde el verano de 2027) \\
\bottomrule
\end{tabular}
\fuente{Elaboración propia a partir de las fuentes citadas en el tema.}
\end{table}

% ==========================================================================
\section{Anexo: la sostenibilidad y el sobreendeudamiento}\label{anexo:deuda}

\begin{itemize}
\item \emph{El saldo estabilizador}: haciendo $b_t=b_{t-1}=b$ en~\eqref{eq:deuda-externa}, $b(1+\gamma_Y)=(1+r^*)\,b+d(1+\gamma_Y)$, de donde $-d\,(1+\gamma_Y)=(r^*-\gamma_Y)\,b$, que es~\eqref{eq:saldo-estabilizador}. \findemo
\item \emph{La solvencia}: despejando $b_{t-1}$ en~\eqref{eq:deuda-externa} y sustituyendo hacia delante, con $r^*>\gamma_Y$ y sin juego de \autor{Ponzi} (la deuda no crece indefinidamente más deprisa que $(1+r^*)/(1+\gamma_Y)$),
  \[
  b_{t-1}=\sum_{k=0}^{\infty}\left(\frac{1+\gamma_Y}{1+r^*}\right)^{k+1}\left(-d_{t+k}\right):
  \]
  la deuda de hoy es el valor actual de los superávits comerciales futuros. \findemo
\item \emph{La deuda en divisas}: con $1+r^*\equiv(1+i^*)/(1+\pi^*)$ el tipo real en la moneda de la deuda y $\gamma_q$ la tasa de depreciación real \vertema{3.B.12},
  \[
  b_t=\frac{(1+r^*)(1+\gamma_q)}{1+\gamma_Y}\,b_{t-1}+d_t ,
  \]
  de modo que una depreciación real ($\gamma_q>0$) sube la pendiente de la recta y acerca la bola de nieve.
\item \emph{El sobreendeudamiento en un modelo sencillo} (letras locales de este anexo):
  \begin{itemize}
  \item \emph{Supuestos}: el país debe $B$, que vence mañana; invierte hoy $I$ y mañana produce $Y\left(I\right)$, con $Y'>0$ e $Y''<0$. Los acreedores pueden obligarle a pagar como mucho una parte $\theta\in(0,1)$ de su producción, con las sanciones de \textcite{BulowRogoff1989JPE}, así que el pago es $\min\{B,\;\theta\,Y\left(I\right)\}$.
  \item \emph{Sin sobreendeudamiento} ($B\leq\theta\,Y\left(I\right)$): el país paga $B$ y se queda con $Y\left(I\right)-B-I$; invierte hasta $Y'\left(I\right)=1$, la inversión eficiente.
  \item \emph{Con sobreendeudamiento} ($B>\theta\,Y\left(I\right)$): paga $\theta\,Y\left(I\right)$ y se queda con $(1-\theta)\,Y\left(I\right)-I$; invierte hasta
    \[
    (1-\theta)\,Y'\left(I\right)=1 ,
    \]
    menos que lo eficiente: \uline{la deuda que no se puede pagar actúa como un impuesto de tipo $\theta$ sobre la inversión}. \findemo
  \item \emph{El tramo decreciente}: con una producción incierta, cuanto mayor es $B$, en más escenarios el país acaba pagando $\theta\,Y$ en lugar de $B$; la cuña afecta a más escenarios y la inversión esperada cae. A partir de cierto $B$, la caída de $Y\left(I\right)$ pesa más que el mayor derecho a cobrar y $V(B)$ baja: es la curva de Laffer de la deuda\nocite{Krugman1988Overhang,Sachs1989Overhang}.
  \item \emph{La quita}: a la derecha del máximo, reducir $B$ quita la cuña en algunos escenarios, sube la inversión y sube $V(B)$. A la izquierda, es un regalo del acreedor al deudor, y una recompra beneficia sobre todo a los acreedores que no venden\nocite{BulowRogoff1988Buyback}.
  \item \emph{El gráfico}: el gráfico~\ref{fig:laffer-deuda} usa $V(B)=B$ hasta $\underline B=2$ y $V(B)=B-(B-2)^2/8$ después, con el máximo en $B^{\max}=6$; la quita de $B_1=8{,}5$ a $B_2=7$ sube el cobro esperado de 3,22 a 3,88.
  \end{itemize}
\end{itemize}

% ==========================================================================
\section{Anexo: umbrales y requisitos}\label{anexo:umbrales}

\begin{table}[H]
\centering
\caption{Umbrales del LIC DSF según la capacidad de endeudamiento}\label{tab:umbrales-licdsf}
\small
\begin{tabular}{>{\raggedright\arraybackslash}m{6.8cm}>{\centering\arraybackslash}m{2.1cm}>{\centering\arraybackslash}m{2.1cm}>{\centering\arraybackslash}m{2.1cm}}
\toprule
\rowcolor{tceeazulpalido} Indicador (deuda externa pública y con garantía pública) & Débil & Media & Fuerte \\
\midrule
Valor actual de la deuda / PIB (\%) & 30 & 40 & 55 \\
Valor actual de la deuda / exportaciones (\%) & 140 & 180 & 240 \\
Servicio de la deuda / exportaciones (\%) & 10 & 15 & 21 \\
Servicio de la deuda / ingresos fiscales (\%) & 14 & 18 & 23 \\
\rowcolor{tceesalvia} Referencia: valor actual de la deuda pública total / PIB (\%) & 35 & 55 & 70 \\
\bottomrule
\end{tabular}
\fuente{Elaboración propia a partir de \textcite{FMI2018LICDSF}. Umbrales vigentes hasta que se aplique el marco revisado en 2026 \parencite{FMI2026LICDSF}.}
\end{table}

\begin{table}[H]
\centering
\caption{La Iniciativa para los PPME reforzada (1999)}\label{tab:ppme}
\small
\begin{tabular}{>{\raggedright\arraybackslash}m{3.6cm}>{\raggedright\arraybackslash}m{10.4cm}}
\toprule
\rowcolor{tceeazulpalido} Elemento & Contenido \\
\midrule
Requisitos & Poder recibir solo préstamos muy favorables (de la AIF y del FMI); una deuda insostenible tras los mecanismos tradicionales de alivio; un historial de reformas con programas del FMI y del Banco Mundial; una estrategia de lucha contra la pobreza elaborada con participación de la sociedad \\
Umbrales de deuda & Valor actual de la deuda externa superior al 150~\% de las exportaciones o, en las economías muy abiertas (exportaciones de al menos el 30~\% del PIB e ingresos fiscales de al menos el 15~\% del PIB), al 250~\% de los ingresos fiscales \\
Punto de decisión & Alivio provisional del servicio de la deuda \\
Punto de culminación & Reducción definitiva de todos los acreedores, a la que se suma la IADM \\
\bottomrule
\end{tabular}
\fuente{Elaboración propia a partir de \textcite{FMI1999HIPCReforzada} y \textcite{FMI2023HIPC}.}
\end{table}

\begin{anotaciones}
\begin{itemize}
\item \emph{Datos para el test}: Club de París, 1956 y 22 miembros permanentes; Plan \autor{Baker}, 1985; Plan \autor{Brady}, 1989; PPME, 1996 y 1999; IADM, 2005; ISSD, de mayo de 2020 a diciembre de 2021; Marco Común, noviembre de 2020; Consenso OCDE, 1978; elemento de donación mínimo, 45, 15 y 10~\%; conferencias de financiación del desarrollo en Monterrey (2002), Doha (2008), Addis Abeba (2015) y Sevilla (2025).
\end{itemize}
\end{anotaciones}

% ==========================================================================
\section{Anexo: el tema en el temario}\label{anexo:temas}

\begin{itemize}
\item \emph{Economía abierta}: la balanza de pagos y la posición de inversión internacional \vertema{3.B.11}; el enfoque intertemporal y la sostenibilidad exterior \vertema{3.B.12}; las crisis de balanza de pagos, las paradas súbitas y el pecado original \vertema{3.B.17}; los movimientos internacionales de capital \vertema{3.B.18}; el sistema económico internacional desde Bretton Woods \vertema{3.B.21}.
\item \emph{Organismos y foros}: el G-20, la OCDE y el CAD \vertema{3.B.19}; el FMI \vertema{3.B.22}; el Banco Mundial y los bancos regionales \vertema{3.B.30}.
\item \emph{Desarrollo}: las teorías del desarrollo y la evaluación experimental \vertema{3.B.28}; el cambio climático y su financiación \vertema{3.B.31}.
\item \emph{Unión Europea y España}: la cooperación al desarrollo de la Unión \vertema{3.B.42}; los instrumentos financieros de apoyo a la internacionalización y la conversión de deuda \vertema{4.A.24}; la cooperación española \vertema{4.A.26}.
\item \emph{Deuda pública}: la misma aritmética, con otras letras \vertema{3.A.39}.
\end{itemize}

% ==========================================================================
\section{Anexo: preguntas de test}\label{anexo:test}

\begin{anotaciones}
Preguntas tipo test del primer ejercicio: \href{https://www.victorgutierrezmarcos.es/oposicion/temario/primer-ejercicio/test/simulador.html?tema=3.B.29}{simulador de test}.
\end{anotaciones}

\bibliografia


\end{document}
